% Exercises for Chapter 3 (Lebesgue measure and measurable functions). Statements only:
% solutions are kept privately. Throughout, m is the Borel-Lebesgue measure on B(R) and
% (X, A, mu) a measure space, as in the notes.

\begin{exercise}[The length is $\sigma$-additive]\label{ex:mf-length}
\hard This exercise proves the two facts of Step~1 that the notes state without proof. Let
$\mathcal E=\{[a,b):0\le a\le b\le1\}$, let $\mathcal A$ be the algebra of finite disjoint
unions of elements of $\mathcal E$, and let $m([a,b))=b-a$. You may use that $\mathcal A$ is an
algebra, and that $[0,1]$ is compact: a family of closed subsets of $[0,1]$ whose intersection
is empty has a finite subfamily whose intersection is empty.
\begin{enumerate}[label=(\alph*)]
  \item Show that if $[a,b)$ is the disjoint union of $[a_1,b_1),\dots,[a_k,b_k)$, then
    $b-a=\sum_{i=1}^k(b_i-a_i)$. Deduce that $m(A)=\sum_im(I_i)$ does not depend on the
    decomposition $A=\bigcup_iI_i$ of $A\in\mathcal A$ into disjoint intervals of $\mathcal E$.
    \emph{Hint:} compare two decompositions with the decomposition into the intervals
    $I_i\cap J_j$.
  \item Show that $m$ is finitely additive on $\mathcal A$, monotone, and finitely
    subadditive: $m(B_1\cup\dots\cup B_k)\le\sum_im(B_i)$ for $B_i\in\mathcal A$.
  \item Let $\mu$ be a finitely additive function on an algebra $\mathcal A$ of subsets of $X$,
    with $\mu(X)<+\infty$. Show that if $\mu(A_n)\to0$ for every decreasing sequence
    $(A_n)_{n\ge1}$ in $\mathcal A$ with $\bigcap_nA_n=\emptyset$, then $\mu$ is
    $\sigma$-additive on $\mathcal A$.
  \item Let $(A_n)_{n\ge1}$ in $\mathcal A$ decrease to $\emptyset$, and let $\varepsilon>0$.
    Show that there are $B_n\in\mathcal A$ whose closures satisfy
    $\overline{B_n}\subseteq A_n$, with $m(A_n\setminus B_n)<\varepsilon2^{-n}$. Show that
    $\bigcap_{n\le N}\overline{B_n}=\emptyset$ for some $N$, deduce that $m(A_N)<\varepsilon$, and
    conclude that $m$ is $\sigma$-additive on $\mathcal A$.
  \item Where does the proof use that the intervals are bounded? Show that on the algebra of
    subsets of $\mathbb N$ that are finite or have a finite complement, the function equal to
    $0$ on finite sets and $1$ on the others is finitely additive but not $\sigma$-additive,
    and say which step of (d) has no analogue there.
\end{enumerate}
\end{exercise}

\begin{exercise}[Steps 2 and 3 of the construction]\label{ex:mf-steps}
\medium Let $Y\subseteq\mathbb R$ be a Borel set. You may use the extension theorem
(\noteref{thm:mf-extension}), the result of Step~1 for every interval $[a,b)$ (that is, the
existence of $m_{[a,b)}$, a measure on $\sigma(\mathcal A_{[a,b)})$ extending the length, where
$\mathcal A_{[a,b)}$ is the algebra of finite disjoint unions of intervals $[x,y)\subseteq[a,b)$),
and that $\mathcal B(\mathbb R)=\sigma(\{[a,b):a<b\})$ (\noteref{prop:sa-borel-generators}).
\begin{enumerate}[label=(\alph*)]
  \item Let $\mathcal C$ be a family of subsets of $\mathbb R$. Show that
    $\{B\cap Y:B\in\sigma(\mathcal C)\}$ is the $\sigma$-algebra on $Y$ generated by
    $\{C\cap Y:C\in\mathcal C\}$. \emph{Hint:} the sets $B\in\sigma(\mathcal C)$ with
    $B\cap Y\in\sigma(\{C\cap Y:C\in\mathcal C\})$ form a $\sigma$-algebra.
  \item Deduce that $\sigma(\mathcal A_{[a,b)})=\{B\in\mathcal B(\mathbb R):B\subseteq[a,b)\}$, and
    that this is also $\mathcal B([a,b))$, the Borel $\sigma$-algebra of the metric space
    $[a,b)$, whose open sets are the sets $U\cap[a,b)$ with $U$ open in $\mathbb R$.
  \item Let $[a,b)\subseteq[c,d)$. Show that $\mathcal B([a,b))\subseteq\mathcal B([c,d))$ and
    that $m_{[a,b)}(A)=m_{[c,d)}(A)$ for every $A\in\mathcal B([a,b))$.
  \item Show that $m(A)=\lim_nm_{[-n,n)}(A\cap[-n,n))$ defines a measure on
    $\mathcal B(\mathbb R)$, and that $m([a,b))=b-a$ for all $a<b$.
  \item Show that $m$ is the only measure on $\mathcal B(\mathbb R)$ with $m([a,b))=b-a$ for all
    $a<b$.
\end{enumerate}
\end{exercise}

\begin{exercise}[Lebesgue sets and completeness]\label{ex:mf-lebesgue-sets}
\medium Let $(X,\mathcal A,\mu)$ be a measure space, with completion
$(X,\mathcal F_\mu(\mathcal A),\bar\mu)$ as in \noteref{def:mf-completion}. You may use
\noteref{prop:mf-completion} and \noteref{prop:mf-limits}.
\begin{enumerate}[label=(\alph*)]
  \item Show that $E\in\mathcal F_\mu(\mathcal A)$ if and only if there are $A,B\in\mathcal A$ with
    $A\subseteq E\subseteq B$ and $\mu(B\setminus A)=0$, and that then $\bar\mu(E)=\mu(A)=\mu(B)$.
  \item Let $\nu$ be a complete measure on a $\sigma$-algebra $\mathcal G\supseteq\mathcal A$
    with $\nu=\mu$ on $\mathcal A$. Show that $\mathcal F_\mu(\mathcal A)\subseteq\mathcal G$ and
    $\nu=\bar\mu$ on $\mathcal F_\mu(\mathcal A)$: the completion is the smallest complete
    extension.
  \item Suppose that $\mu$ is complete. Show that if $f:X\to\overline{\mathbb R}$ is measurable
    and $g:X\to\overline{\mathbb R}$ satisfies $\{f\ne g\}\subseteq N$ for some negligible $N$,
    then $g$ is measurable; and that if $f_n\to f$ a.e.\ with every $f_n$ measurable, then $f$
    is measurable.
  \item In the counterexample \noteref{exa:mf-ae-counterexample}, show that on
    $(\mathbb R,\mathcal L(\mathbb R),\bar m)$ the function $g=\1_A$ is measurable and
    $f=g$ a.e.
\end{enumerate}
\end{exercise}

\begin{exercise}[A subset of the Cantor set that is not Borel]\label{ex:mf-cantor-nonborel}
\hard This exercise proves the claim of \noteref{rem:mf-cantor}. Let $C$ be the Cantor set and
$c:[0,1]\to[0,1]$ the Cantor function: on $C$,
$c\big(\sum_jd_j3^{-j}\big)=\sum_j\frac{d_j}22^{-j}$ ($d_j\in\{0,2\}$), and $c$ is constant on
each of the open intervals that make up $[0,1]\setminus C$. You may use that $c$ is continuous,
non-decreasing and onto $[0,1]$, that $[0,1]\setminus C$ is a countable disjoint union of open
intervals, and \emph{Vitali's theorem}: every $E\in\mathcal L(\mathbb R)$ with $\bar m(E)>0$ has a
subset that is not in $\mathcal L(\mathbb R)$.
\begin{enumerate}[label=(\alph*)]
  \item Show that $h(x)=x+c(x)$ is a strictly increasing continuous bijection of $[0,1]$ onto
    $[0,2]$, with a continuous inverse.
  \item Show that $h$ maps each interval of $[0,1]\setminus C$ onto an interval of the same
    length. Deduce that $m(h([0,1]\setminus C))=1$ and $m(h(C))=1$.
  \item Let $V\subseteq h(C)$ be a set that is not Lebesgue measurable, and
    $A=h^{-1}(V)\subseteq C$. Show that $A\in\mathcal L(\mathbb R)$ but $A\notin\mathcal B(\mathbb R)$.
  \item Conclude that $\mathcal B(\mathbb R)$ is not complete with respect to $m$, and that a
    homeomorphism of $\mathbb R$ can map a Lebesgue measurable set onto a set that is not
    Lebesgue measurable.
\end{enumerate}
\end{exercise}

\begin{exercise}[The $\sigma$-algebra generated by a function]\label{ex:mf-sigma-f}
\easy This is the exercise left to the reader after \noteref{def:mf-generated}. Let
$f:X\to(Y,\mathcal E)$.
\begin{enumerate}[label=(\alph*)]
  \item Show that $f^{-1}(\mathcal E)=\{f^{-1}(B):B\in\mathcal E\}$ is a $\sigma$-algebra on $X$,
    and deduce that $\sigma(f)=f^{-1}(\mathcal E)$ (\noteref{prop:mf-sigma-f}).
  \item Show that $f$ is measurable from $(X,\mathcal A)$ to $(Y,\mathcal E)$ if and only if
    $\sigma(f)\subseteq\mathcal A$.
  \item Show that if $\mathcal E=\sigma(\mathcal G)$, then
    $\sigma(f)=\sigma(\{f^{-1}(G):G\in\mathcal G\})$.
  \item Let $X=Y=\mathbb R$ with Borel sets and $f(x)=x^2$. Show that $\sigma(f)$ consists of
    the Borel sets $S$ that are symmetric, $S=-S$. Is $x\mapsto x$ measurable with respect to
    $\sigma(f)$? Is $x\mapsto|x|$?
\end{enumerate}
\end{exercise}

\begin{exercise}[Suprema, limits, and the extended real line]\label{ex:mf-limits}
\medium These are the proofs that the notes leave as exercises
(\noteref{prop:mf-limits} and \noteref{cor:mf-pointwise-limit}). Functions take values in
$\overline{\mathbb R}$ and are measurable in the sense of \noteref{rem:mf-extended}. You may use
\noteref{prop:mf-generators} and \noteref{prop:sa-borel-generators}.
\begin{enumerate}[label=(\alph*)]
  \item Show that $g:X\to\overline{\mathbb R}$ is measurable if and only if $\{g>t\}\in\mathcal A$
    for every $t\in\mathbb R$, if and only if $\{g\le t\}\in\mathcal A$ for every $t\in\mathbb R$.
  \item Let $f_n:X\to\overline{\mathbb R}$, $n\ge1$, be measurable. Show that $\sup_nf_n$
    and $\inf_nf_n$ are measurable, then that $\liminf_nf_n$ and $\limsup_nf_n$ are.
  \item Deduce \noteref{cor:mf-pointwise-limit}.
  \item Let $f,g:X\to\overline{\mathbb R}$ be measurable. Show that $\{f<g\}$, $\{f\ne g\}$ and
    $\{f=g\}$ are in $\mathcal A$. Deduce that the set of points where $(f_n(x))_n$ converges
    in $\overline{\mathbb R}$ is in $\mathcal A$.
\end{enumerate}
\end{exercise}

\begin{exercise}[$\sigma(f)$ on a die]\label{ex:mf-finite-sigma}
\easy This is a computational exercise. Let $X=\{1,\dots,6\}$, the faces of a die, with
$\mathcal A=\mathcal P(X)$, and on it
\[
  f(x)=\1_{\{x\text{ even}\}},\qquad g(x)=|x-3.5|,\qquad h(x)=\lceil x/2\rceil,\qquad
  k(x)=\1_{\{x\ge5\}},
\]
each with values in a finite set $Y$ equipped with $\mathcal P(Y)$.
\begin{enumerate}[label=(\alph*)]
  \item Write a function that returns $\sigma(f)=\{f^{-1}(B):B\subseteq f(X)\}$ for a function
    on a finite set, and compute the atoms (the minimal non-empty sets) and the number of
    elements of $\sigma(f)$, $\sigma(g)$, $\sigma(h)$, $\sigma(k)$. Check on random functions
    that $\sigma(f)$ is a $\sigma$-algebra with $2^{|f(X)|}$ elements.
  \item Using the preimage criterion, decide which of $f,g,h,k$ are $\sigma(g)$-measurable and
    which are $\sigma(h)$-measurable. For each $\sigma(h)$-measurable one, find $\varphi$ with
    $u=\varphi\circ h$.
  \item Compute $\sigma(g)\cap\sigma(h)$, $\sigma(g)\cup\sigma(h)$ and $\sigma(g,h)$. Which of
    them are $\sigma$-algebras?
  \item Let $F=(0,1,2,2,1,0)$ on $\{0,\dots,5\}$, with values in $Y=\{0,1,2\}$, and
    $\mathcal A'=\sigma(\{\{0,5\}\})$. Check that $F^{-1}(\{0\})\in\mathcal A'$ but that $F$ is not
    $\mathcal A'$-measurable. Why does this not contradict \noteref{prop:mf-generators}?
\end{enumerate}
\end{exercise}

\begin{exercise}[Cantor sets, thin and fat]\label{ex:mf-cantor-sets}
\medium This is a computational exercise. For $0<r\le\frac13$, let $F_0=[0,1]$ and obtain
$F_j$ from $F_{j-1}$ by removing, from the middle of each of its $2^{j-1}$ closed intervals,
an open interval of length $r^j$. For $r=\frac13$, $F_j=C_j$, the stages of the Cantor set.
\begin{enumerate}[label=(\alph*)]
  \item Write vectorised code that returns the intervals of $F_k$ as a $2^k\times2$ array and
    computes $m(F_k)$. Check that $m(C_k)=(\frac23)^k$ for $k\le20$, and find, using the formula
    (not by building $C_k$), the first $k$ with $m(C_k)<10^{-6}$.
  \item Let $L_j$ be the length of the intervals of $F_j$. Show that $L_j=(L_{j-1}-r^j)/2$, that
    every removal fits ($L_{j-1}>r^j$) when $r\le\frac13$, and that $F=\bigcap_jF_j$ has
    measure $\frac{1-3r}{1-2r}$. Which $r$ gives $m(F)=\frac12$, and which gives $m(F)=0.99$?
    Check numerically.
  \item For $r=\frac14$, find $L_j$ in closed form and show that $F$ contains no interval,
    although $m(F)=\frac12$.
  \item Draw $10^4$ random points of $C$ by choosing $40$ ternary digits $d_j\in\{0,2\}$
    independently and uniformly. Check that they lie in $C_{20}$, and that their images under the
    Cantor function look uniformly distributed on $[0,1]$ (a Kolmogorov--Smirnov test). Why does
    this illustrate that $C$ is uncountable?
\end{enumerate}
\end{exercise}

\begin{exercise}[Monte Carlo with an honest error bar]\label{ex:mf-monte-carlo}
\medium This is a computational exercise. Let $A=\{x\in(0,1]:\sin(1/x)>0\}$.
\begin{enumerate}[label=(\alph*)]
  \item Show that $A$ is a countable disjoint union of intervals, hence a Borel set, and that
    $m(A)=1-\frac{\ln2}{\pi}\approx0.77936$. \emph{Hint:} $\sum_{k\ge1}(-1)^{k+1}/k=\ln2$.
  \item Let $U_1,\dots,U_N$ be independent and uniform on $(0,1)$, and
    $\hat p_N=\frac1N\sum_i\1_A(U_i)$. Compute $\hat p_N$ for $N=10^3,\dots,10^6$, with the
    asymptotic $95\%$ interval $\hat p_N\pm1.96\sqrt{\hat p_N(1-\hat p_N)/N}$ and the
    \emph{Hoeffding} interval $\hat p_N\pm\sqrt{\ln(2/0.05)/(2N)}$, which has probability at
    least $95\%$ of covering $m(A)$ for every $N$, by Hoeffding's inequality
    $\mathbb P(|\hat p_N-p|\ge\varepsilon)\le2e^{-2N\varepsilon^2}$.
  \item Repeat the experiment $200$ times with $N=10^4$ and report how often each interval
    covers $m(A)$. Which interval is wider, and by how much?
  \item How many samples does the Hoeffding bound require to guarantee an error below
    $10^{-3}$ with probability $99\%$?
  \item The same code would estimate $m(\mathbb Q\cap(0,1))$ as $1$, whatever $N$. Why? What does
    this say about which sets a sampler can ``see''?
\end{enumerate}
\end{exercise}

\begin{exercise}[How fast do the simple functions converge?]\label{ex:mf-dyadic-rate}
\medium This is a computational exercise on the functions $f_n$ of \noteref{eq:mf-dyadic}.
\begin{enumerate}[label=(\alph*)]
  \item Implement $f_n=\min(n,2^{-n}\lfloor2^nf\rfloor)$ on arrays of values of $f$, and check
    it against the case-by-case definition \noteref{eq:mf-dyadic}.
  \item Let $f(x)=12x(1-x)$ on $[0,1]$. Compute $\sup_x(f(x)-f_n(x))$ on a fine grid for
    $n=1,\dots,14$, and explain the values: what happens for $n=1,2$, and why is the supremum
    exactly $2^{-n}$ for $n\ge3$?
  \item Let $f(x)=-\log x$ on $(0,1)$. Show that $\sup_x(f(x)-f_n(x))=+\infty$ for every $n$, and
    that the (improper Riemann) integral of the error is
    \[
      E_n=\int_0^1\big(f(x)-f_n(x)\big)\,\mathrm dx
      =e^{-n}+\big(1-e^{-n}\big)\frac{1-e^{-h}(1+h)}{1-e^{-h}},\qquad h=2^{-n}.
    \]
    \emph{Hint:} substitute $x=e^{-y}$, and treat $\{f\ge n\}$ and each
    $\{\frac{i-1}{2^n}\le f<\frac{i}{2^n}\}$ separately.
  \item Check the formula numerically for a few $n$, and show that $E_n\sim2^{-n-1}$. From
    which $n$ on does the rounding error dominate the error due to the cut at height $n$?
\end{enumerate}
\end{exercise}
